\section*{ЗАКЛЮЧЕНИЕ}
\addcontentsline{toc}{section}{ЗАКЛЮЧЕНИЕ}

В ходе выполнения данной работы была создана программно-информационная система для автоматизации деятельности учебно-методического центра повышения квалификации. Система включает в себя веб-интерфейс для работы с данными о слушателях и десктоп-приложение для управления встроенными серверами PostgreSQL и Flask.

Разработанная система позволяет централизованно хранить информацию о слушателях, программах обучения, группах и выдаваемых документах, обеспечивает гибкую фильтрацию и поиск по множеству критериев, формирование статистических отчётов, импорт данных из Excel-файлов, а также разграничение прав доступа (администратор, обычный пользователь).

Основные результаты работы:

\begin{enumerate}
	\item Проведён анализ предметной области и существующих систем автоматизации учебного процесса. Выявлены недостатки традиционного подхода и обоснована необходимость разработки собственной информационной системы.
	\item Разработана концептуальная модель системы, построена ER-диаграмма и реляционная модель базы данных. Определены функциональные и нефункциональные требования.
	\item Осуществлено проектирование архитектуры программной системы (клиент-сервер с веб-интерфейсом и десктоп-приложением). Разработаны структура базы данных, пользовательский интерфейс и программный интерфейс приложения (API).
	\item Реализована и протестирована программно-информационная система. Проведено системное тестирование в соответствии со сценариями использования, подтвердившее работоспособность всех функций.
\end{enumerate}

Все требования, объявленные в техническом задании, были полностью реализованы, все задачи, поставленные в начале разработки проекта, решены. Разработанная система готова к внедрению в учебно-методическом центре повышения квалификации.  
